
\chapter*{Abstract}

This thesis evaluates the CarenR distributional subgroup discovery framework as a rule-based machine learning approach for two complementary tasks on real-world agricultural production time series: interpretable pattern extraction (Goal 1) and walk-forward prediction (Goal 2). The application domain is annual wine production forecasting for two contrasting Portuguese wine regions, the Douro Demarcated Region (RDD, 89 years, 1934--2022) and the Vinho Verde Region (RVV, 80 years, 1942--2021). The two regions serve as a natural comparative experiment, same pipeline, same algorithms, different structural histories and climate regimes, to evaluate when and why the method succeeds or fails under different signal conditions.

Two methodological contributions support the dual goals. For Goal 1, CarenR is applied to the full linearly detrended production series to extract interpretable agroclimatic rules. Each rule is a conjunction of climate conditions under which the distribution of production residuals shifts significantly, assessed by a two-sample Kolmogorov--Smirnov test. Rule quality is further checked using LOESS-detrended residuals: features selected by CarenR retain 90--92\% of their correlation after aggressive era-trend removal, confirming genuine within-era climate signal rather than era-membership artefacts. For Goal 2, a two-decision decomposition framework separates prediction error into trend quality (Decision 1: how well the training trend extrapolates to test years) and climate signal extraction (Decision 2: whether a rule-based model can predict detrended residuals better than the training median), identifying which element of the pipeline is the bottleneck. The evaluation uses an expanding-window walk-forward protocol with 18 scenarios for RDD and 16 for RVV, with all preprocessing --- detrending, feature discretisation, and feature selection --- performed within each training fold to prevent leakage. Eleven CarenR variants are evaluated as a structured sensitivity analysis of feature selection strategies, alongside RIPPER, M5Rules, a Decision Tree regressor, and naive baselines, totalling 17 models.

Goal 1 (Rule Discovery) is achieved\footnote{\textbf{[Defesa] P:} O título diz ``forecasting'' mas a previsão falha --- é honesto? \textbf{R:} A previsão é o objeto de estudo, não uma promessa. O achado central é a assimetria discovery vs.\ prediction; o resultado nulo é um diagnóstico, não um fracasso.}. CarenR's KS-test-based subgroup induction discovers 48 statistically validated rules for RDD and 20 for RVV. Feature frequency analysis shows that harvest-period soil water appears in 54\% of RDD rules, the most consistently selected variable across all 48 rules, while in RVV the most frequently selected features are heat-stress days around veraison (up to 50\% of rules) and harvest soil water (25\%), with the Leaf Area Index (LAI) family present in \erratamark{4 of the 20 rules}{era ``all 20''}. LOESS detrending validates that the selected features carry genuine within-era signal (correlations drop by ~8\% after era-trend removal), confirming that CarenR extracts real structure rather than era-membership artefacts. A cross-algorithm comparison shows that RIPPER and M5Rules independently converge on the same variable families, confirming that the identified variables are robust across three different optimisation criteria.

Goal 2 (Predictive Validation) is not achieved overall\footnote{\textbf{[Defesa] P:} Se a correlação existe (r cerca de 0,47), porque não prevê? \textbf{R:} A discovery só precisa que o sinal exista; a prediction precisa que ele domine o erro out-of-sample. Com CV cerca de 47\%, o clima explica ~22\% da variância dos resíduos --- os restantes ~78\% são noise.}. No rule-based model consistently outperforms the Naive\_Median baseline across the full evaluation sequence. The best CarenR variant falls 3\% behind the baseline in RDD (189 vs 184 mhl weighted MAE; Wilcoxon $p = 0.468$ across all 18 scenarios, not significant) and 22\% behind in RVV (172 vs 140 mhl). In the 9 RVV scenarios where rules fired, CarenR was significantly worse than the baseline (Wilcoxon $p = 0.004$). A post-hoc era-composition analysis, conducted after observing the results and therefore exploratory, suggests CarenR consistently outperforms the baseline in RDD when training data reaches ~29\% post-1990 representation (5 of 5 consecutive scenarios, binomial $p = 0.031$), pointing to a conditional regime where rule-based prediction could succeed. The two-decision framework explains the failure: the demands of accurate trend modelling and sufficient residual signal for rule induction are in fundamental tension under a time-only detrending covariate.

\textbf{Keywords: }\textit{distributional subgroup discovery, CarenR, rule-based machine learning, walk-forward validation, time series evaluation, two-decision decomposition, feature selection, wine production forecasting, non-stationarity, era confound}

\cleardoublepage

\chapter*{Resumo}

Esta tese avalia o método de descoberta de subgrupos distributivos CarenR enquanto abordagem de aprendizagem automática baseada em regras para duas tarefas complementares em problemas de séries temporais de produção agrícola do mundo real: extração de padrões interpretáveis (Objetivo~1) e previsão com validação \textit{walk-forward} (Objetivo~2). O domínio de aplicação corresponde à previsão anual da produção de vinho em duas regiões vitivinícolas portuguesas contrastantes: a Região Demarcada do Douro (RDD, 89 anos, 1934--2022) e a Região dos Vinhos Verdes (RVV, 80 anos, 1942--2021). Estas duas regiões constituem uma experiência comparativa natural, na qual o mesmo \textit{pipeline} e os mesmos algoritmos são aplicados a históricos estruturais e regimes climáticos distintos, permitindo avaliar em que condições e por que razão o método obtém sucesso ou falha perante diferentes regimes de sinal.

Duas contribuições metodológicas sustentam estes dois objetivos. Para o Objetivo~1, o CarenR é aplicado à série temporal completa da produção, após remoção da tendência linear, com o propósito de extrair regras agroclimáticas interpretáveis. Cada regra consiste numa conjunção de condições climáticas sob as quais a distribuição dos resíduos da produção difere significativamente, sendo essa diferença avaliada através de um teste de Kolmogorov--Smirnov de duas amostras. A qualidade de cada regra é posteriormente validada recorrendo aos resíduos obtidos após a remoção da tendência através do método LOESS, demonstrando que as características selecionadas pelo CarenR preservam entre 90\% e 92\% da sua correlação, mesmo após uma remoção agressiva da tendência associada à era, o que confirma a existência de um sinal climático genuíno intra-era, em vez de artefactos decorrentes da pertença à era.

Para o Objetivo~2, é proposta uma estrutura de decomposição em duas decisões, que separa o erro de previsão em duas componentes: qualidade da modelação da tendência (Decisão 1: até que ponto a tendência estimada no conjunto de treino é capaz de extrapolar para os anos de teste) e extração do sinal climático (Decisão 2: se um modelo baseado em regras consegue prever os resíduos obtidos após a remoção da tendência melhor do que a mediana do conjunto de treino). Esta decomposição permite identificar qual das componentes do \textit{pipeline} constitui o principal fator limitativo do desempenho preditivo.

A avaliação da capacidade preditiva segue um protocolo de validação \textit{walk-forward} com janela de treino em expansão, compreendendo 18 cenários para a RDD e 16 para a RVV. Todo o pré-processamento --- remoção da tendência, discretização e seleção de atributos --- é realizado exclusivamente dentro de cada conjunto de treino, de forma a evitar a fuga de informação (\textit{data leakage}). São avaliadas onze variantes do CarenR, correspondentes a uma análise de sensibilidade das estratégias de seleção de atributos, em conjunto com os algoritmos RIPPER e M5Rules, um regressor baseado em Árvores de Decisão e modelos de base simples, perfazendo um total de 17 modelos.

O Objetivo~1 (Descoberta de Regras) é alcançado. A indução de subgrupos baseada no teste de Kolmogorov--Smirnov implementada pelo CarenR identifica 48 regras estatisticamente significativas para a RDD e 20 para a RVV. A análise da frequência dos atributos revela que a água do solo durante o período de colheita está presente em 54\% das regras da RDD, sendo a variável selecionada de forma mais consistente ao longo das 48 regras, enquanto o Índice de Área Foliar (LAI) na colheita domina o conjunto de regras da RVV, estando presente em todas as 20 regras.

A validação através da remoção da tendência pelo método LOESS confirma que os atributos selecionados transportam um sinal genuíno intra-era, verificando-se apenas uma redução de 8\% nas correlações após a remoção da tendência da era, o que demonstra que o CarenR extrai uma estrutura climática real, em vez de artefactos associados à dinâmica temporal. Adicionalmente, a comparação entre algoritmos evidencia que o RIPPER e o M5Rules convergem, de forma independente, para as mesmas famílias de variáveis, reforçando a robustez das variáveis identificadas sob três critérios distintos de otimização.

O Objetivo~2 (Validação Preditiva), contudo, não é alcançado de forma geral. Nenhum dos modelos baseados em regras supera consistentemente a linha de base Naive\_Median (mediana simples) ao longo de toda a sequência de avaliação. A melhor variante do CarenR apresenta um desempenho 3\% inferior ao da linha de base na RDD (MAE ponderado de 189 versus 184 mhl; teste de Wilcoxon ($p = 0{,}468$) para os 18 cenários, sem diferença estatisticamente significativa) e 22\% inferior na RVV (172 versus 140 mhl). Nos nove cenários da RVV em que as regras são ativadas, o CarenR apresenta um desempenho significativamente inferior ao da linha de base (teste de Wilcoxon, $p = 0{,}004$).

Uma análise post hoc da composição temporal das diferentes eras, realizada após a observação dos resultados e, por conseguinte, de natureza exploratória, sugere que o CarenR supera consistentemente a linha de base na RDD quando o conjunto de treino atinge uma representação de pelo menos 29\% de observações posteriores a 1990 (5 em 5 cenários consecutivos; teste binomial, $p = 0{,}031$), apontando para a existência de um regime condicional no qual a previsão baseada em regras poderá ser eficaz.

Por fim, esta estrutura de decomposição em duas decisões permite explicar o insucesso preditivo observado: as exigências simultâneas de uma modelação robusta da tendência e da preservação de um sinal residual suficiente para a indução de regras encontram-se numa tensão fundamental quando a remoção da tendência depende exclusivamente de uma covariável temporal.

\textbf{Palavras-chave: }\textit{descoberta de subgrupos distributivos, CarenR, aprendizagem automática baseada em regras, validação walk-forward, avaliação de séries temporais, seleção de atributos, previsão da produção de vinho, não estacionariedade}
